\documentclass[10pt,a4paper]{amsart}
\usepackage[margin=19mm]{geometry}
\usepackage{amsmath,amssymb,mathtools}
\usepackage[scr=rsfs,cal=euler]{mathalfa}
\usepackage{xfrac}
\usepackage[unicode,hidelinks,bookmarks=false]{hyperref}
\newcommand{\ie}{i.e.\ }
\newcommand{\cf}{cf.\ }
\newcommand{\resp}{respectively\ }
\newcommand{\Subsection}{\subsection}
\providecommand{\nobreakdash}{\nobreak\hbox{-}}
\setlength{\emergencystretch}{2em}
\multlinegap0pt
\usepackage{amssymb}
\usepackage{mathtools}
\usepackage{xcolor}
\usepackage{dsfont}
\usepackage{float}
\usepackage[shortlabels]{enumitem}
\newtheorem{theorem}{Theorem}
\title{Two-Orbit Polytopes}
\author{Isabel Hubard, Egon Schulte}
\date{Source: arXiv:2604.00185v1 (CC BY 4.0)}
\begin{document}
\maketitle

\noindent\emph{Source and licence.} Two-Orbit Polytopes. Version arXiv:2604.00185v1 (CC BY 4.0), distributed by arXiv under the Creative Commons Attribution 4.0 International licence, http://creativecommons.org/licenses/by/4.0/, as recorded in the arXiv licence metadata for this exact version and shown on the abstract page https://arxiv.org/abs/2604.00185v1. Full licence notice: \texttt{licenses/arxiv-2604.00185-CC-BY-4.0.txt}.\par\smallskip

\noindent\textit{Editorial scope.} Counted unit: the article's theorem generators (source0.tex lines 526--529). The article's complete original proof of it is delivered with it (source0.tex lines 531--600, one \texttt{proof} environment of 70 lines). No mathematics is written, repaired or re-derived: the statement and the proof are the article's own lines, in the article's own notation, and no retained line is substituted or re-flowed.  Retained uncounted as the source-local context the proof needs: lines 504--511.  Nothing was omitted from any retained span, so the package carries no omission record.  No retained line contains a citation, so the wrapper carries no bibliography and no bibliography entry.  No retained line contains a figure environment, an image-inclusion command or drawing code, so no image is delivered and the package declares no asset dependency.  Editorial formatting only, in addition: an AMS \texttt{amsart} preamble that restates the article's own declarations and macros used by the retained lines, namely the environments \texttt{theorem} on separate counters, together with the standard packages those lines need.  The retained environments print in this document as Theorem 1 in order of appearance, whereas the article prints the same items with the numbering of its own full document.  The complete native source of the article as distributed in the arXiv e-print for this exact version is bundled unchanged as \texttt{sources/197572/source0.tex}.
Our first theorem says that  
\begin{eqnarray}
\label{gi}
\mathcal{G}_{I} :=\mathcal{G}_{I}(\Phi):= 
\{ \rho_i \,|\, i \in I\} \,\cup\, \{\alpha_{j,k}\,|\, j,k \notin I   \} \,\cup\, \{ \alpha_{j,i,j} \,|\, i \in I, j \notin I \}
\end{eqnarray}
is a set of generators of $\Gamma(\mathcal{P})$. We usually suppress the reference to $\Phi$ in the notation if no confusion is possible.
\smallskip

\begin{theorem}
\label{generators}
Let $\mathcal{P}$ be a two-orbit polytope in the class $2_I$, with $I\subset N$, and let $\mathcal{G}_{I}$ be as in (\ref{gi}). Then $\Gamma(\mathcal{P})$ is generated by $\mathcal{G}_{I}$.
\end{theorem}

\begin{proof}
As above, let $\Phi$ be the base flag of $\mathcal{P}$. Our goal is to write an arbitrary element $\psi$ of $\Gamma(\mathcal{P})$ in terms of the elements of $\mathcal{G}_{I}$. 

Consider the flag $\Psi:=\Phi\psi$, which belongs to the same flag orbit as $\Phi$. By the strong flag connectedness of~$\mathcal{P}$, there exists a sequence of successively adjacent flags 
\begin{eqnarray}
\label{adjacentsequence1}
\Phi, \Phi^{i_1}, \Phi^{i_1, i_2}, \dots, \Phi^{i_1, i_2, \dots, i_l}=\Psi,
\end{eqnarray}
all containing $\Phi \cap \Psi$, joining $\Phi$ and $\Psi$.  For the purpose of this proof define $\Phi^{(0)}:=\Phi$ and $\Phi^{(m)}:= \Phi^{i_1,\dots, i_m}$ for $m=1,\ldots,l$. Our goal is to exploit the structure of the sequence in (\ref{adjacentsequence1}) to produce a factorization of $\psi$ into elements from $\mathcal{G}_{I}$.
 
Since $\Phi$ and $\Psi$ are in the same orbit, the number of superscripts $i_k$ that are not contained in~$I$ must be even, equal to $2s$ (say). Let $i_{k_1}, \ldots, i_{k_{2s}}$ denote the terms in the sequence $i_1, i_2, \ldots, i_l$ that are not contained in~$I$, and set $k_0:=0$ and $k_{2s+1}:=l+1$. Then~(\ref{adjacentsequence1}) takes the form
\begin{eqnarray}
\label{seq2}
\Phi=\Phi^{(k_0)},  \Phi^{(1)}, \dots,  \Phi^{(k_1)},\ldots, \Phi^{(k_2)}, \ldots\ldots, \Phi^{(k_{2s-1})}, \ldots, \Phi^{(k_{2s})}, \ldots, \Phi^{(l)} = \Psi.
\end{eqnarray}
Here the flags $\Phi^{(1)}, \ldots, \Phi^{(k_1-1)}$ at the beginning of the sequence are in the same orbit as~$\Phi$, since all of the original superscripts involved lie in $I$. However, starting with $q=1$, each time we encounter a term of the form $\Phi^{(k_q)}$ as we move along the sequence in (\ref{seq2}), we change from one flag orbit to the other flag orbit. On the other hand, no change of flag orbit occurs at the other terms in (\ref{seq2}).

Now for $q=1, \dots, 2s-1$ define 
\begin{eqnarray*}
\beta_q&\!:=\!&\alpha_{i_{k_q},i_{k_{q+1}}} \alpha_{i_{k_q}, i_{k_{q+1}-1}, i_{k_q}} \alpha_{i_{k_q}, i_{k_{q+1}-2}, i_{k_q}} \ldots\ldots \alpha_{i_{k_q}, i_{k_q+1} i_{k_q}}, \\
\gamma_q&\!:=\!& \rho_{i_{k_{q+1}-1}} \rho_{i_{k_{q+1}-2}} \ldots\ldots \rho_{i_{k_q+1}};
\end{eqnarray*}
here, if $k_q$ and $k_{q+1}$ are consecutive integers then $\beta_{q}=\alpha_{i_{k_q},i_{k_{q+1}}}$ and $\gamma_{q}=1$. We also set
\begin{eqnarray*}
\gamma_{0}   &\!:=\!& \rho_{i_{k_{1}-1}} \rho_{i_{k_{1}-2}} \ldots \rho_{i_1},\\
\gamma_{2s} &\!:=\!& \rho_{i_l} \rho_{i_{l-1}} \dots \rho_{i_{k_{2s}+1}} ,
\end{eqnarray*}
so in particular, $\gamma_{0}=1$ if $k_{1}=1$, and $\gamma_{2s}=1$ if $k_{2s}=l$.
Note that for each $q$, each factor occurring in the expressions for $\beta_q$ or $\gamma_q$ is an element of $\mathcal{G}_{I}$. 

Next we show that the action of the generators of $\mathcal{G}_{I}$ on $\Phi$ gives
\begin{eqnarray*}
\Phi \beta_q      &\!=\!& \Phi^{i_{k_{q}}, i_{k_q+1}, \dots, i_{k_{q+1}-1} , i_{k_{q+1}}} \;\,(q=1,\ldots,2s-1), \\ 
\Phi \gamma_q &\!=\!& \Phi^{i_{k_q+1}, i_{k_q+2}, \dots, i_{k_{q+1}-1}} \;\,(q=0,\ldots,2s).
\end{eqnarray*}
Here the right hand side of the second equation is to be read as $\Phi$ if either $q=0$ and $k_{1}=1$, or $q=2s$ and $k_{2s}=l$. Then the equation for $\Phi\beta_q $ can be established as follows:
\[\begin{array}{lll}
\Phi \beta_q 
&=& \Phi \alpha_{i_{k_q},i_{k_{q+1}}} \alpha_{i_{k_q}, i_{k_{q+1}-1}, i_{k_q}} \alpha_{i_{k_q}, i_{k_{q+1}-2}, i_{k_q}} \ldots\ldots \alpha_{i_{k_q}, i_{k_q+1} i_{k_q}}\\[.05in]
&=& (\Phi^{i_{k_q},i_{k_{q+1}}}) \alpha_{i_{k_q}, i_{k_{q+1}-1}, i_{k_q}} \alpha_{i_{k_q}, i_{k_{q+1}-2}, i_{k_q}} \ldots\ldots \alpha_{i_{k_q}, i_{k_q+1} i_{k_q}}\\[.05in]
&=& (\Phi \alpha_{i_{k_q}, i_{k_{q+1}-1}, i_{k_q}} \alpha_{i_{k_q}, i_{k_{q+1}-2}, i_{k_q}} \ldots\ldots \alpha_{i_{k_q}, i_{k_q+1} i_{k_q}})^{i_{k_q},i_{k_{q+1}}}\\[.05in]
&=& ((\Phi^{i_{k_q}, i_{k_{q+1}-1}, i_{k_q}}) \alpha_{i_{k_q}, i_{k_{q+1}-2}, i_{k_q}} \ldots\ldots \alpha_{i_{k_q}, i_{k_q+1} i_{k_q}})^{i_{k_q},i_{k_{q+1}}}\\[.05in]
&=& ((\Phi \alpha_{i_{k_q}, i_{k_{q+1}-2}, i_{k_q}} \ldots\ldots \alpha_{i_{k_q}, i_{k_q+1} i_{k_q}})^{i_{k_q}, i_{k_{q+1}-1}, i_{k_q}})^{i_{k_q},i_{k_{q+1}}}\\[.05in]
&=& (\Phi \alpha_{i_{k_q}, i_{k_{q+1}-2}, i_{k_q}} \ldots\ldots \alpha_{i_{k_q}, i_{k_q+1} i_{k_q}})^{i_{k_q}, i_{k_{q+1}-1},i_{k_{q+1}}}\\[.05in]
&=& (\Phi \alpha_{i_{k_q}, i_{k_{q+1}-3}, i_{k_q}} \ldots\ldots \alpha_{i_{k_q}, i_{k_q+1} i_{k_q}})^{i_{k_q},i_{k_{q+1}-2},i_{k_{q+1}-1},i_{k_{q+1}}}\\[.02in]
&\vdots&\\[.02in]
&=& \Phi^{i_{k_{q}}, i_{k_q+1}, \dots, i_{k_{q+1}-1} , i_{k_{q+1}}}.
\end{array} \]
The equation for $\Phi \gamma_q$ is more straightforward and follows similarly.

Now define the element $\widehat{\psi}$ of the subgroup $\langle\mathcal{G}_{I}\rangle$ of $\Gamma(\mathcal{P})$ by
\begin{eqnarray}
\label{star}
\widehat{\psi}: =\gamma_{2s} \beta_{2s-1} \gamma_{2s-2} \beta_{2s-3} \ldots\ldots \gamma_4 \beta_3 \gamma_2 \beta_1 \gamma_{0}.
\end{eqnarray}
Then the image of $\Phi$ under $\widehat{\psi}$ is given by
\[ \begin{array}{lll}
\Phi\widehat{\psi} 
&=& (\Phi\gamma_{2s})\,\beta_{2s-1}\gamma_{s-2}\beta_{2s-2}\gamma_{2s-3}\ldots\ldots\beta_{1}\gamma_{0}\\[.05in]
&=& (\Phi^{\,i_{k_{2s}+1},i_{k_{2s}+2},\ldots,i_l})\,\beta_{2s-1}\gamma_{s-2}\beta_{2s-2}\gamma_{2s-3}\ldots\ldots\beta_{1}\gamma_{0}\\[.05in]
&=& (\Phi\beta_{2s-1}\gamma_{s-2}\beta_{2s-2}\gamma_{2s-3}\ldots\ldots\beta_{1}\gamma_{0})^{\,i_{k_{2s}+1},i_{k_{2s}+2},\ldots,i_l}\\[.05in]
&=& ((\Phi\beta_{2s-1})\gamma_{s-2}\beta_{2s-2}\gamma_{2s-3}\ldots\ldots\beta_{1}\gamma_{0})^{\,i_{k_{2s}+1},i_{k_{2s}+2},\ldots,i_l}\\[.05in]
&=& ((\Phi^{i_{k_{2s-1}}, i_{k_{2s-1}+1}, \dots, i_{k_{2s}-1} , i_{k_{2s}}})\gamma_{s-2}\beta_{2s-2}\gamma_{2s-3}\ldots\ldots\beta_{1}\gamma_{0})^{\,i_{k_{2s}+1},i_{k_{2s}+2},\ldots,i_l}\\[.05in]
&=& (\Phi\gamma_{s-2}\beta_{2s-2}\gamma_{2s-3}\ldots\ldots\beta_{1}\gamma_{0})^{\,{i_{k_{2s-1}}, i_{k_{2s-1}+1}, \ldots, i_{k_{2s}-1},i_{k_{2s}}},i_{k_{2s}+1},i_{k_{2s}+2},\ldots,i_l}\\[.02in]
&\vdots&\\[.02in]
&=&\, \Phi^{i_1,\ldots,i_l} \,=\, \Psi \,=\, \Phi\psi.
\end{array} \]

Therefore, since the images of $\Phi$ under $\psi$ and $\widehat {\psi}$ coincide, we must have $\psi=\widehat{\psi}\in\langle\mathcal{G}_{I}\rangle$. This completes the proof.
\end{proof}

\end{document}
